\documentclass[11pt]{article}
\usepackage[margin=1in]{geometry}
\usepackage{fontspec}
\setmainfont{Latin Modern Roman}[SmallCapsFont={Latin Modern Roman Caps}]
\usepackage{amsmath,amssymb}
\usepackage{xcolor}
\usepackage{enumitem}
\usepackage{array}
\usepackage{longtable}
\usepackage{booktabs}
\usepackage{hyperref}
\usepackage[normalem]{ulem}
\hypersetup{colorlinks=true, linkcolor=blue, urlcolor=blue}

\newcommand{\response}[1]{\smallskip\noindent{\color{blue!60!black}\textbf{Trả lời.}}\ #1\par\medskip}
\newcommand{\loc}[1]{\textcolor{blue!50!black}{[#1]}}
% original English wording of a quotation, shown for checking
\newcommand{\en}[1]{{\color{gray}\small(EN: ``#1'')}}

\newenvironment{point}[1]{\medskip\noindent\textbf{\uline{#1}}\par}{\par}

\title{\vspace{-2em}Thư trả lời phản biện\\[0.3em]
\large Bản thảo: \emph{Thuật toán Chiếu Thích nghi Hai Track cho một bài toán chấp nhận tách phi tuyến với các ràng buộc hỗn hợp lồi--giả lồi}\\[0.2em]
\normalsize Tạp chí: \emph{Optimization Methods and Software}}
\author{Nguyen Van The, Tran Ngoc Thang và Nguyen Hong Son}
\date{}

\begin{document}
\maketitle
\vspace{-2em}

\begin{center}
\fbox{\parbox{0.92\textwidth}{\small Bản dịch tiếng Việt của \texttt{DTAP Finally v2/response\_letter.tex}, chỉ dùng để rà soát nội dung, không nộp. Ý kiến của Reviewer cũng được dịch. Với các câu trích từ bản thảo, nguyên văn tiếng Anh được ghi kèm bằng chữ xám ``(EN: \dots)'' để tiện đối chiếu.}}
\end{center}

\noindent Chúng tôi cảm ơn Editor và cả hai Reviewer đã đọc kỹ và góp ý mang tính xây dựng.
Các ý kiến đã làm bài báo vững hơn đáng kể, và ở nhiều chỗ đã khiến
chúng tôi bỏ đi các khẳng định quá mức và bổ sung bằng chứng mới, trung thực hơn. Dưới đây chúng tôi trả lời
từng ý kiến và chỉ ra thay đổi tương ứng trong bản thảo sửa.
\textbf{Mọi phần chữ được thêm hoặc sửa trong bản thảo sửa đều được in bằng
\textcolor{red}{chữ màu}} để dễ tìm thấy các thay đổi. Số mục, công thức,
định lý và bảng dưới đây theo bản thảo sửa. Tóm lại, bản sửa:
\begin{itemize}[nosep]
  \item gộp phần Giới thiệu và phần Các công trình liên quan cũ thành một phần Giới thiệu
        duy nhất, ngắn hơn, và rút bài từ $39$ xuống $25$ trang, viết lại toàn bộ
        bằng các câu ngắn hơn;
  \item làm rõ giới hạn của khẳng định ``không cần hằng số'' (Nhận xét~2 mới) và trình bày lại trọng số tỉ lệ
        như một cách phân bổ kinh nghiệm theo tỉ phần phần dư, không khẳng định là cần thiết hay
        tối ưu (Nhận xét~1 mới);
  \item làm tường minh lập luận hội tụ của toàn bộ dãy (viết lại Giai đoạn~3 của
        Định lý~3.12), làm rõ rằng nó chỉ cần một điểm tụ khả thi và
        tính đơn điệu Fej\'er, không cần giới hạn duy nhất;
  \item bổ sung một phần chẩn đoán bộ giải (Mục~4.4 mới) báo cáo mức hoạt động của ràng buộc
        theo từng quỹ đạo, một phép so sánh với phương pháp dưới gradient có mô-đun tinh chỉnh cùng chi phí tinh chỉnh, và
        các bài toán có kiểm soát với một ràng buộc giả lồi thực sự không lồi;
  \item chứng minh rằng tỉ số dùng trong chọn đặc trưng thực ra là lồi (Mệnh đề~4.1 mới),
        nói rõ rằng ứng dụng là một trường hợp riêng lồi, và bỏ mọi
        khẳng định về ưu thế đồng đều;
  \item thống nhất định dạng tài liệu tham khảo, bổ sung các công trình về bước tự thích nghi và phân rã
        cho SFP theo yêu cầu, chuyển các kết quả ở phụ lục vào thân bài, và giải quyết
        mọi vấn đề về ký hiệu và trình bày đã được nêu.
\end{itemize}

% =====================================================================
\section*{Trả lời Editor}

\begin{point}{Associate Editor không ghi ý kiến cụ thể nào; quyết định là sửa lớn (major
revision).}
\response{Chúng tôi đã sửa bản thảo theo đúng báo cáo của cả hai Reviewer,
như trình bày chi tiết dưới đây, và đã đánh dấu mọi thay đổi bằng chữ màu. Hai thay đổi
cấu trúc lớn nhất --- gộp và rút ngắn phần mở đầu (Reviewer~2)
và bổ sung phần chẩn đoán bộ giải với bằng chứng về mức hoạt động của ràng buộc và phương pháp cơ sở đã tinh chỉnh
(Reviewer~1) --- đã được tóm tắt trong danh sách ở trên.}
\end{point}

% =====================================================================
\section*{Trả lời Reviewer \#1}

\noindent Chúng tôi cảm ơn Reviewer đã tóm tắt chính xác đóng góp của bài và đã đưa ra năm
ý kiến trực tiếp giúp phát biểu lý thuyết cẩn thận hơn và phần thực nghiệm
đầy đủ hơn.

\begin{point}{R1.1 Trọng số mềm $\mu_k=J_k^{\mathrm{cvx}}/(J_k^{\mathrm{cvx}}+J_k^{\mathrm{ps}})$
là trung tâm của tính mới của thuật toán, nhưng chứng minh hội tụ chỉ đòi $\mu_k$
thỏa một số điều kiện khả tổng -- bất kỳ dãy nào trong $[0,1]$ cũng đủ. Vì vậy dạng
được chọn không cần thiết về mặt cấu trúc. Tác giả nên hoặc biện minh cho cách đặt trọng số
cụ thể này, hoặc trình bày lại đóng góp cho phù hợp.}
\response{Chúng tôi chọn đề nghị thứ hai của Reviewer và trình bày lại đóng góp.
Chúng tôi đồng ý rằng dạng tỉ lệ không cần thiết về mặt cấu trúc: bất đẳng thức
tựa Fej\'er một bước (Bổ đề~3.8) đúng với mọi $\mu_k\in[0,1]$, và chứng minh
hội tụ chỉ cần trọng số giữ cho mỗi track còn hoạt động --- chẳng hạn với trọng số hằng
$\mu_k\equiv\tfrac12$, chứng minh của Bổ đề~3.10 và của Giai đoạn~2 vẫn đúng với
vài thay đổi nhỏ. Thực nghiệm cũng cho thấy điều tương tự: một biến thể cân bằng, gán
trọng số bằng nhau khi cả hai track bị vi phạm, thành công ở $26/36$ lần chạy trên các bài toán
đặc trưng và $24/36$ trên họ bài toán không lồi, đúng bằng số lần của DTAP \loc{Mục~4.4,
Bảng~7}. Vì vậy chúng tôi không còn trình bày trọng số cụ thể này như một điều cần thiết về lý thuyết hay
một lựa chọn tối ưu. Phần Giới thiệu nay nói rằng ``chứng minh chỉ dùng trọng số này để giữ cho một
track bị vi phạm kéo dài vẫn có trọng số dương'' \en{the proof uses this weight only to keep a
persistently violated track active}. Nhận xét~1 mới \loc{Nhận xét~1, Mục~3.4} nói rõ điều này: nó ghi lại chỗ
dạng tỉ lệ được dùng trong chứng minh hiện tại và lưu ý rằng bất đẳng thức một bước,
dù đúng với mọi $\mu_k\in[0,1]$, tự nó không suy ra hội tụ với mọi
dãy trọng số. Các chế độ suy biến của Hệ quả~3.13--3.14 ứng với
$\mu_k\equiv1$ và $\mu_k\equiv0$. Cấu trúc hai track
--- một bước co theo phần dư ở nơi có bất đẳng thức gradient định lượng,
và một bước có chặn an toàn ở nơi chỉ có thông tin về dấu --- vẫn là
đóng góp; còn trọng số cụ thể dùng để pha trộn chúng thì không.}
\end{point}

\begin{point}{R1.2 Khẳng định ``không cần hằng số'' bị phóng đại. Thuật toán không
đòi tường minh hằng số Lipschitz, nhưng Giả thiết~1(A2) và Bổ đề~4.9 bảo đảm
dãy lặp nằm trong một tập compact, nơi chắc chắn tồn tại một hằng số Lipschitz. Khẳng định
này cần được nói rõ giới hạn.}
\response{Chúng tôi đồng ý. Sự phụ thuộc ngầm như vậy có ở mọi phương pháp có dãy lặp bị chặn
và gradient liên tục, nên nay chúng tôi trình bày vấn đề theo nghĩa \emph{đầu vào} của thuật toán
chứ không theo nghĩa sự tồn tại của một hằng số. Phát biểu chính xác hiện nay là: không có
chuẩn toán tử, hằng số Lipschitz hay mô-đun H\"older nào là \emph{đầu vào} của
phép lặp, trong khi \emph{phân tích} có dùng các cận hữu hạn của giá trị hàm và
gradient trên tập compact chứa dãy lặp (đã được chứng minh là bị chặn). Nhận xét~2
mới \loc{Nhận xét~2, Mục~3.4} nói rõ điều này: ``Gradient bị chặn kéo theo các hàm
Lipschitz trên tập compact lồi đó, nhưng các cận này không phải là đầu vào của
thuật toán'' \en{Bounded gradients imply
Lipschitz continuity of the functions on that convex compact set, but these bounds are
not algorithm inputs}, và lưu ý rằng tính bị chặn được chứng minh trước (Bổ đề~3.9, tức
Bổ đề~4.9 cũ) rồi mới dùng đến các cận như vậy. Tóm tắt nay nói rằng
độ dài bước ``không cần biết trước chuẩn toán tử, hằng số Lipschitz hay mô-đun H\"older''
\en{require no prescribed operator norm, Lipschitz constant, or H\"older modulus}
và rằng ``chứng minh dùng các cận địa phương trên miền bị chặn chứa dãy lặp'' \en{the proof uses local bounds on the bounded iterate region}, còn phần Kết luận nói
rằng ``không có mô-đun chính quy nào là đầu vào, mặc dù phân tích dùng các cận trên một tập compact
chứa dãy lặp'' \en{no regularity modulus is an input, although the analysis uses bounds on a compact
set containing the iterates}. Chúng tôi không còn dùng cụm ``không cần hằng số'' mà không kèm giới hạn.}
\end{point}

\begin{point}{R1.3 Định lý~4.12 khẳng định toàn bộ dãy hội tụ tới một điểm khả
thi. Tuy nhiên, Giai đoạn~2 của chứng minh chỉ xác lập được rằng có ít nhất một điểm giới hạn
khả thi. Vì bài toán không lồi, có thể tồn tại nhiều điểm giới hạn. Tác giả
nên hoặc làm chặt lập luận, hoặc sửa định lý thành hội tụ theo
dãy con.}
\response{Chúng tôi giữ phát biểu về hội tụ của toàn bộ dãy, vì nó đúng, và viết lại
Giai đoạn~3 để bước mà Reviewer chỉ ra nay được trình bày tường minh. Lập luận không
cần điểm giới hạn duy nhất. Với \emph{mọi} $x^\star\in\Omega$ cố định, Bổ đề~3.9(i)
cho thấy $\{\lVert x^k-x^\star\rVert^2\}$ hội tụ (tính đơn điệu Fej\'er đối với
toàn bộ tập khả thi, chỉ dùng $R^{x^\star}_k\ge0$ từ Bổ đề~3.5). Giai đoạn~2
tạo ra \emph{một} điểm tụ khả thi $\bar x\in C\cap\Omega$. Áp dụng Bổ đề~3.9(i)
với lựa chọn hợp lệ $x^\star=\bar x$, giới hạn $\lim_k\lVert x^k-\bar x\rVert^2$
tồn tại; vì có một dãy con hội tụ tới $\bar x$, giới hạn đó bằng $0$, nên \emph{toàn bộ}
dãy hội tụ tới $\bar x$. Do đó không thể có nhiều điểm tụ khả thi cùng tồn tại.
Giai đoạn~3 viết lại trình bày đầy đủ điều này và kết thúc bằng câu: ``Lập luận không
cần $\Omega$ chỉ có một điểm: một điểm tụ khả thi và sự hội tụ của khoảng cách tới nó
là đủ'' \en{The argument does not
require $\Omega$ to be a singleton: one feasible cluster point and the convergence of its
distance suffice} \loc{chứng minh Định lý~3.12,
Mục~3.8}. Chúng tôi cũng sửa phần Giới thiệu (đóng góp~C2) cho khớp với cách diễn đạt này.}
\end{point}

\begin{point}{R1.4 Các ngưỡng $b_i$ được đặt bằng các phân vị của $F_i(w)$ lấy mẫu để
bảo đảm $\Omega\neq\emptyset$. Trong ứng dụng thực tế, ngưỡng do bài toán cho trước,
không điều chỉnh được. Quan trọng hơn, nếu $b_3$ lớn đến mức $F_3(w)\ll b_3$ ở mọi nơi, thì
ràng buộc giả lồi không bao giờ chặt. Tác giả nên báo cáo tần suất ràng buộc hoạt động
và xác nhận rằng bối cảnh hỗn hợp thực sự được kích hoạt.}
\response{Chúng tôi nay trả lời trực tiếp cả hai nửa của ý kiến này, trong tiểu mục mới
\emph{Mức hoạt động của ràng buộc và một phép so sánh với phương pháp dưới gradient đã tinh chỉnh} \loc{Mục~4.4}. Về
ngưỡng: chúng tôi làm rõ rằng ngưỡng được cho từ bên ngoài vẫn là đầu vào cố định, và rằng
việc hiệu chỉnh theo phân vị chỉ được dùng để xác định bài toán chuẩn, với cùng
các ngưỡng dùng chung cho mọi bộ giải; tính nhất quán ($\Omega\neq\emptyset$) được xem là một
giả thiết, và các mức cụ thể được ghi trong siêu dữ liệu bổ sung
\loc{Mục~4.1.3 và Mục~4.4}. Về mức hoạt động: chúng tôi báo cáo tần suất
$F_3(w^k)-b_3>10^{-6}$ dọc theo mọi quỹ đạo. Ràng buộc thứ ba thực ra
\emph{không hoạt động} trên bảy trong tám bộ dữ liệu và chỉ hoạt động tại $1$ trong $48$ điểm được
đánh giá trên Abalone (Bảng~6). Chúng tôi nói thẳng hệ quả, đồng ý với
Reviewer: trên các bài toán này track giả lồi hầu như không bao giờ chặt, nên
cơ chế hai track không cần thiết ở đó --- các bài toán chọn đặc trưng này chỉ kiểm tra một
lớp con lồi. Thật vậy, Mệnh đề~4.1 mới chứng minh rằng
tiêu chí phân thức $F_3=w^{\top}Q_Rw/\rho^{\top}w$ là \emph{lồi} trên một lân cận
của đơn hình khi $Q_R\succeq0$ và $\min_k\rho_k>0$, nên mô hình đặc trưng là
một NSFP với hộp chuẩn mà mọi ràng buộc đều lồi. Để thực sự kích hoạt bối cảnh hỗn hợp, chúng tôi bổ sung các bài toán
có kiểm soát với một ràng buộc giả lồi được chứng minh là không lồi,
$F_3(x)=\kappa\log(1+\lVert x-e_1\rVert^2)$ \loc{Công thức~(47), Mục~4.4}; ở đó
track giả lồi hoạt động trong cả $36$ lần chạy (mức hoạt động gộp $12\,835/12\,877$ điểm
được đánh giá), nên track giả lồi được kiểm tra trực tiếp chứ không phải một cách hình thức.}
\end{point}

\begin{point}{R1.5 DTAP và SG-const/SG-dyn không phân biệt được về mặt thống kê trên ba
trong bốn độ đo (Bảng~5). Thất bại của SG-dyn trên $3$ bộ dữ liệu được quy cho việc ước lượng
mô-đun H\"older kém --- nhưng với một mô-đun đã tinh chỉnh, nó có thể cho kết quả tương đương. Một
thực nghiệm bổ sung với mô-đun tinh chỉnh thủ công sẽ làm phép so sánh vững hơn.}
\response{Chúng tôi đã bổ sung đúng thực nghiệm này \loc{Mục~4.4, Bảng~6--8}. Với SG-dyn, bên cạnh
biến thể dùng mô-đun lấy mẫu, chúng tôi nay báo cáo một biến thể \emph{tinh chỉnh} có mô-đun được
chọn từ một lưới thủ công $c\in\{10^{-3},\dots,10^{1}\}$ trên các điểm xuất phát tách riêng, và
tính riêng chi phí tinh chỉnh ($179\,744$ lần đánh giá trong tối ưu và
$10.80$\,s trên các bài toán đặc trưng; số liệu tương tự cho họ bài toán tổng hợp). Tinh chỉnh
nâng số lần thành công từ điểm xuất phát chuẩn của SG-dyn từ $5/9$ lên $9/9$ và, như Reviewer
dự đoán, một phương pháp dưới gradient đã tinh chỉnh có thể bằng hoặc tốt hơn DTAP --- chẳng hạn trên Abalone,
và trên họ bài toán không lồi, SG-const tinh chỉnh và SG-dyn tinh chỉnh đạt sai số dừng ở cả
$36$ lần chạy, so với $24/36$ của DTAP (Bảng~7). Vì vậy chúng tôi đã sửa mọi khẳng định
so sánh: bài nay nói rằng không phép so sánh từng cặp nào có ý nghĩa sau hiệu chỉnh
Holm, rằng kết quả phân loại phản ánh sự đánh đổi giữa kích thước và khả năng dự đoán chứ không phải
ưu thế, và, trong phần Kết luận, rằng ``các phương pháp dưới gradient đã tinh chỉnh có thể bằng hoặc tốt hơn
DTAP'' \en{tuned subgradient methods can match or outperform
DTAP}. Ưu điểm được báo cáo của DTAP được thu hẹp lại ở chỗ không có mô-đun nào phải ước lượng hay
tinh chỉnh như một đầu vào của phép lặp.}
\end{point}

% =====================================================================
\section*{Trả lời Reviewer \#2}

\noindent Chúng tôi cảm ơn Reviewer đã đọc chi tiết từng trang. Chúng tôi đã làm theo
mọi ý kiến; bài nay ngắn hơn đáng kể và, chúng tôi hy vọng, rõ ràng hơn nhiều.

\subsection*{Các vấn đề chính}

\begin{point}{M1. Mục~1 (Giới thiệu) và Mục~2 (Các công trình liên quan) nên được sắp xếp lại và
gộp vào Mục~1, giới hạn khoảng $4$ trang; bài nên được rút ngắn (khoảng
$25$ trang) bằng cách bỏ nội dung thừa và các câu dài.}
\response{Đã làm. Mục Các công trình liên quan (và Bảng~1 của nó) đã được bỏ, và
nội dung cốt lõi được cô đọng vào phần Giới thiệu \loc{Mục~1}, nay nhắc lại
phương pháp $CQ$, các chiến lược bước tự thích nghi và phân rã, và các công trình phi tuyến
gần nhất (Li và cộng sự, Yen và cộng sự, Hu và cộng sự) trong vài đoạn ngắn, và nêu ba
đóng góp thay vì năm. Phần Kiến thức chuẩn bị được rút gọn chỉ còn các kết quả thực sự được dùng
và không còn chia tiểu mục; khung chọn đặc trưng (Mục~5 cũ) nay là
một tiểu mục của phần thực nghiệm \loc{Mục~4.1}; biến thể thưa DTAP-S, phép so sánh
thời gian chạy, mục về độ bền/ablation, và phần ``Tóm tắt kết quả'' đã được bỏ;
phần Kết luận chỉ còn một đoạn. Bài nay dài $25$ trang, giảm từ $39$.}
\end{point}

\begin{point}{M2. Cần cải thiện cách viết đáng kể: nhiều câu
mơ hồ hoặc quá dài. Hãy mô tả bài toán, thuật toán, v.v.\ bằng các câu ngắn; trau chuốt
từng câu; bỏ các mô tả ``vô nghĩa''.}
\response{Chúng tôi đã viết lại bản thảo theo yêu cầu này. Các câu dài được tách ra,
các cụm từ mơ hồ được thay hoặc xóa, và các đoạn mô tả không phục vụ
kết quả nào đã được bỏ. Tóm tắt, phần Giới thiệu, phần Kiến thức chuẩn bị, phần mở đầu của Mục~3,
và toàn bộ phần thực nghiệm đã được viết lại. Tên trong ngoặc của các bổ đề,
mệnh đề, nhận xét và hệ quả đã được bỏ để rút ngắn văn bản.}
\end{point}

\begin{point}{M3. Định dạng tài liệu tham khảo không nhất quán, và thiếu các tài liệu quan trọng về
SFP; ví dụ, tác giả của [23] và [35] cũng có các công trình về bước tự thích nghi hoặc phân rã
cho SFP.}
\response{Chúng tôi đã thống nhất định dạng tài liệu tham khảo theo một kiểu (tên viết tắt, tên tạp chí
in nghiêng, tập, số, trang, năm, và DOI nếu có), và nay trích dẫn các công trình
về bước tự thích nghi và phân rã của tác giả [23] (Z.~Li, D.~Han, W.~Zhang) và
[35] (Q.~Yang) cho SFP (nhiều tập): Zhang--Han--Li (\emph{Inverse Problems}, 2009),
Zhang--Han--Yuan (\emph{Comput.\ Optim.\ Appl.}, 2012), Zhao--Yang (\emph{Inverse
Problems}, 2011), và Yang--Yang--Zhang (\emph{J.\ Optim.\ Theory Appl.}, 2014). Chúng được
trích ở đoạn đầu của phần Giới thiệu, nơi nhắc lại các bước tự thích nghi và
chiến lược phân rã cho SFP. Các tài liệu chỉ cần cho phần nội dung đã bỏ
đã được xóa.}
\end{point}

\begin{point}{M4. Tại sao một số kết quả số được trình bày ở Phụ lục?}
\response{Chúng tôi đồng ý, và phụ lục đã được bỏ. Các giá trị độ chính xác và AUC theo từng bộ dữ liệu
nay nằm trong thân bài \loc{Bảng~3, Mục~4.3}, và hai bộ giải cơ sở được
mô tả trong thân bài \loc{Mục~4.2.1}. Các bảng ngưỡng và thời gian chạy đã được bỏ;
ngưỡng của mọi bài toán được ghi trong gói tái lập
bổ sung.}
\end{point}

\subsection*{Các vấn đề nhỏ}

\noindent{\color{blue!60!black}\textbf{Trả lời.}}\ Chúng tôi đồng ý với cả hai mươi ý kiến
nhỏ và đã làm theo tất cả, không có ngoại lệ. Để ngắn gọn, chúng tôi liệt kê
chúng, cùng thay đổi tương ứng và vị trí, trong bảng dưới đây; trong mọi trường hợp
thay đổi được đánh dấu bằng chữ màu trong bản thảo sửa.

\medskip
{\footnotesize
\setlength{\tabcolsep}{5pt}
\renewcommand{\arraystretch}{1.25}
\begin{longtable}{@{}>{\raggedright\arraybackslash}p{0.42\textwidth}
                    >{\raggedright\arraybackslash}p{0.52\textwidth}@{}}
\toprule
\textbf{Ý kiến nhỏ của Reviewer} & \textbf{Thay đổi đã làm (vị trí trong bản thảo sửa)}\\
\midrule
\endfirsthead
\toprule
\textbf{Ý kiến nhỏ của Reviewer} & \textbf{Thay đổi đã làm (vị trí trong bản thảo sửa)}\\
\midrule
\endhead
\bottomrule
\endfoot
m1 (tr.~2, d.~18) --- cụm ``reference point'' (điểm tham chiếu) mơ hồ.
  & Đã bỏ cụm này; quy tắc của L\'opez và cộng sự nay được trích trong nhóm các phương pháp tự thích nghi ở Mục~1.\\
m2 (tr.~2, d.~27) --- ``differing'' hay ``different densities'' (mật độ khác nhau).
  & Dùng ``different''; câu về xạ trị đã được bỏ (Mục~1).\\
m3 (tr.~2, d.~35) --- trau chuốt câu ``The NSFP is not itself new'' (bản thân NSFP không mới).
  & Đã bỏ câu này; Li và cộng sự, Yen và cộng sự, Hu và cộng sự nay được so sánh trực tiếp với DTAP (Mục~1).\\
m4 (tr.~2, d.~39) --- ``active index set'' hay ``active set'' (tập hoạt động).
  & Dùng ``active set''; câu đó đã bỏ; chúng tôi dùng ``approximate active set'' (tập hoạt động xấp xỉ) (Mục~3.2).\\
m5 (tr.~2, d.~47--49) --- câu ``We ask instead \dots'' (thay vào đó chúng tôi hỏi \dots) không rõ.
  & Đã bỏ; quan hệ giữa các lớp hàm được nêu lại bằng các câu ngắn (Mục~1).\\
m6 (tr.~3, d.~48--53) --- các câu quá dài.
  & Đã bỏ đoạn đó; rút còn một câu (Mục~1).\\
m7 (tr.~3) --- năm đóng góp quá chi tiết.
  & Thay bằng ba mục ngắn C1--C3 (Mục~1).\\
m8 (tr.~5) --- ngữ pháp của câu ``Rule (3) has since been extended'' (quy tắc (3) từ đó đã được mở rộng).
  & Câu này đã bỏ cùng với mục Các công trình liên quan.\\
m9 (tr.~5, d.~33) --- ``drives their proof is for free''; ``nonconvexity sits in the sets''.
  & Đoạn không rõ nghĩa này đã bỏ cùng với mục Các công trình liên quan.\\
m10 (tr.~5, d.~35) --- ``plain gradient inequality'' (bất đẳng thức gradient đơn giản).
  & Đã bỏ (cụm này muốn nói bất đẳng thức gradient của hàm lồi $f(y)\ge f(x)+\langle\nabla f(x),y-x\rangle$).\\
m11 (tr.~5, d.~44) --- ``the closest predecessor'' (công trình tiền nhiệm gần nhất).
  & Đã bỏ; ba công trình đi trước được so sánh trực tiếp với DTAP (Mục~1).\\
m12 (tr.~5, d.~48) --- ``step'' hay ``step size'' (độ dài bước).
  & Dùng ``step size''; câu đó đã bỏ cùng với mục Các công trình liên quan.\\
m13 (tr.~7, d.~31) --- ``dtap'' phải là ``DTAP''.
  & Đã sửa trong toàn bài; từ viết tắt nay luôn là ``DTAP''.\\
m14 (tr.~7, d.~39) --- ``Combettes 10'' phải là ``Combettes~[10]''.
  & Đã sửa; trích dẫn còn lại có dạng ``in the sense of Combettes~[5]'' (theo nghĩa của Combettes [5]) (Mục~3.7).\\
m15 (tr.~9) --- các tiểu mục 3.1--3.2 không cần thiết; 3.3 nên chuyển sang phần thực nghiệm.
  & Phần Kiến thức chuẩn bị không còn chia tiểu mục (Mục~2); mô hình chọn đặc trưng nay là Mục~4.1.\\
m16 (tr.~12) --- $H:=\mathbb{R}^n$ bị lặp lại.
  & Đã bỏ phần lặp khỏi phát biểu bài toán (Mục~3.1).\\
m17 (tr.~13) --- nghĩa của ``standing assumption'' (giả thiết thường trực).
  & Tiểu mục được đổi tên thành ``Assumptions'' (Các giả thiết) (Mục~3.3); tên cũ đã bỏ.\\
m18 (tr.~18, d.~40) --- kiểu AMS cho $\sum_{k=0}^{\infty}\lVert x^{k+1}-x^k\rVert^2<\infty$.
  & Chuỗi được hiển thị với cận trên/dưới (Mệnh đề~3.11); các công thức hiển thị khác đã được kiểm tra.\\
m19 (tr.~20, d.~36) --- trau chuốt cụm ``so no Young splitting is needed'' (nên không cần tách theo Young).
  & Đã bỏ cụm này khỏi chứng minh Hệ quả~3.13.\\
m20 (tr.~20, d.~35) --- biểu thức $K\cap\bigcap_{i\in I^{\mathrm{cvx}}}$ không tốt.
  & Viết lại thành $\Omega=\{x\in K:F_i(x)\le b_i\ \forall i\in I^{\mathrm{cvx}}\}$ (Hệ quả~3.13), và thành $\bigcap_{i\in I^{\mathrm{cvx}}}\{x\in K:F_i(x)\le b_i\}$ trong chứng minh Định lý~3.12.\\
\end{longtable}}

\begin{point}{Chung. Xin hãy xem xét các ý kiến trên và làm cho bài súc tích.}
\response{Chúng tôi đã làm theo mọi ý kiến và rút bài từ $39$ xuống $25$ trang
(xem M1 và M2).}
\end{point}

% =====================================================================
\section*{Các thay đổi khác do tác giả tự thực hiện}

\noindent Trong khi sửa bài, chúng tôi còn thực hiện các thay đổi sau, không do
Reviewer yêu cầu, được liệt kê để đầy đủ và minh bạch.
\begin{enumerate}[leftmargin=1.6em,itemsep=3pt]
\item \textbf{Tính lồi của tiêu chí phân thức.} Bản nộp ban đầu gọi
  $F_3(w)=w^{\top}Q_Rw/\rho^{\top}w$ là giả lồi. Nay chúng tôi chứng minh nó \emph{lồi} trên một
  lân cận mở của đơn hình khi $Q_R\succeq0$ và $\min_k\rho_k>0$
  (Mệnh đề~4.1), nên bài toán chọn đặc trưng là một NSFP mà mọi ràng buộc đều lồi. Cách
  gán $F_3$ vào track giả lồi dùng trong thực nghiệm vẫn hợp lệ
  (lồi $\Rightarrow$ giả lồi), và cách gán toàn lồi được báo cáo như một phương án
  so sánh.
\item \textbf{Miền xác định của các hàm thành phần.} Các hàm nay được định nghĩa trên một lân cận
  mở $U$ của $K$, và (A2) đòi khả vi liên tục trên $U$; điều này
  bao gồm các hàm phân thức có mẫu số dương trên $K$. Một đoạn ngắn sau
  Giả thiết~1 giải thích điều này và lý do $\Omega$ lồi.
\item \textbf{Phạm vi của phân tích.} Ngoài Nhận xét~1 và~2 (xem R1.1, R1.2), Nhận xét~5
  nay nói chính xác rằng cận $O(N^{-1/3})$ chỉ liên quan đến độ vi phạm \emph{tốt nhất}
  của track lồi --- không phải cận cho điểm lặp cuối, không phải tốc độ cho toàn bộ phần dư hay
  cho khoảng cách tới tập khả thi.
\item \textbf{Quy trình thực nghiệm.} Mục~4.2.2 nay nêu chính xác quy trình:
  chuẩn hóa min--max được khớp một lần trên tập huấn luyện, các mặt nạ bỏ phiếu được tính một lần
  rồi dùng lại, việc hiệu chỉnh xác suất của SVM dùng một seed cố định, và việc đánh giá là
  khám phá chứ không phải kiểm định lồng. Chỉ số ổn định Nogueira (với hệ số hiệu chỉnh phương sai đúng
  $N_f/(N_f-1)$) và quy trình thống kê (kiểm định hạng có dấu chính xác,
  hiệu chỉnh Holm, kiểm định Friedman) được nêu tường minh.
\item \textbf{Các giá trị được tính lại.} Các giá trị độ ổn định, AUC của SVM, và
  giá trị $p$ đã được tính lại và nay gồm cả giá trị $p$ gốc và sau hiệu chỉnh Holm; các cột cho
  SG-const và SG-dyn đã được thêm vào. Độ chính xác và mọi kết quả $5$-NN không đổi.
\item \textbf{Các khẳng định được sửa lại.} Các phát biểu về ưu thế (độ chính xác, tốc độ, độ bền)
  đã được bỏ hoặc giới hạn lại: không phép so sánh từng cặp nào có ý nghĩa sau hiệu chỉnh Holm,
  kết quả cho thấy sự đánh đổi giữa kích thước và khả năng dự đoán, và các phương pháp dưới gradient đã tinh chỉnh có thể bằng hoặc
  tốt hơn DTAP.
\item \textbf{Trình bày thuật toán.} Phép kiểm tra dừng của Thuật toán~1 nay được
  đặt ở đầu Bước~1 thay vì ở một Bước~3 riêng; các điểm lặp và
  điểm được trả về không đổi.
\item \textbf{Chặn an toàn ở mẫu số.} Mục~4.1.3 nay báo cáo rằng điểm Fisher
  nhỏ nhất trên mọi bài toán là $4.77\times10^{-5}$, nên chặn an toàn $\max\{\rho^{\top}w,10^{-15}\}$
  trong mã cài đặt không bao giờ có tác dụng và $F_3$ được cài đặt đúng là~(41).
\item \textbf{Các tuyên bố.} Tuyên bố về việc dùng AI tạo sinh và tuyên bố về tính sẵn có của
  dữ liệu đã được cập nhật, và một gói tái lập nay được cung cấp kèm bài.
\end{enumerate}

\bigskip
\noindent Chúng tôi tin rằng các sửa đổi này đã trả lời đầy đủ ý kiến của các Reviewer và Editor
và đã cải thiện độ rõ ràng, chặt chẽ và trung thực của bản thảo. Chúng tôi một lần nữa cảm ơn
các Reviewer và Editor đã dành thời gian.

\end{document}
